\[ %%%%%%%%%%%%%%%%%%%%%%%%%%%%%%%%%%%%%%%%%%%%%%%%%%%%%%%%%%%%%%%%%%%%%%%%%%%%%%%% \section{Data and Code Availability} \label{sec:data_code} %%%%%%%%%%%%%%%%%%%%%%%%%%%%%%%%%%%%%%%%%%%%%%%%%%%%%%%%%%%%%%%%%%%%%%%%%%%%%%%% To support reproducibility and facilitate further research in physics-guided machine learning, the source code, experimental configurations, and generated datasets associated with PHYSGEN will be publicly released. The complete implementation will include: \begin{itemize} \item PHYSGEN model architecture; \item physics-guided graph learning modules; \item stability control components; \item training and evaluation pipelines; \item benchmark generation scripts; \item visualization tools. \end{itemize} %%%%%%%%%%%%%%%%%%%%%%%%%%%%%%%%%%%%%%%%%%%%%%%%%%%%%%%%%%%%%%%%%%%%%%%%%%%%%%%% \subsection{Code Availability} %%%%%%%%%%%%%%%%%%%%%%%%%%%%%%%%%%%%%%%%%%%%%%%%%%%%%%%%%%%%%%%%%%%%%%%%%%%%%%%% The source code of PHYSGEN will be released through a public repository: \begin{center} \texttt{https://github.com/[repository-name]} \end{center} The repository will contain: \begin{itemize} \item complete model implementation; \item configuration files; \item training scripts; \item evaluation scripts; \item experiment reproduction notebooks; \item documentation. \end{itemize} The released code will allow researchers to reproduce all experiments reported in this manuscript, including: \begin{itemize} \item Lorenz-63 forecasting; \item Reaction-Diffusion simulations; \item Navier--Stokes experiments; \item ablation studies; \item stability analysis. \end{itemize} %%%%%%%%%%%%%%%%%%%%%%%%%%%%%%%%%%%%%%%%%%%%%%%%%%%%%%%%%%%%%%%%%%%%%%%%%%%%%%%% \subsection{Dataset Availability} %%%%%%%%%%%%%%%%%%%%%%%%%%%%%%%%%%%%%%%%%%%%%%%%%%%%%%%%%%%%%%%%%%%%%%%%%%%%%%%% The benchmark datasets used in this work are generated using deterministic simulation procedures based on the corresponding governing equations. The generated datasets will be deposited in an open research repository: \begin{center} \texttt{https://zenodo.org/[dataset-record]} \end{center} The dataset release will include: \begin{itemize} \item initial conditions; \item simulation trajectories; \item graph representations; \item physical parameters; \item train/validation/test splits. \end{itemize} %%%%%%%%%%%%%%%%%%%%%%%%%%%%%%%%%%%%%%%%%%%%%%%%%%%%%%%%%%%%%%%%%%%%%%%%%%%%%%%% \subsection{Reproducibility} %%%%%%%%%%%%%%%%%%%%%%%%%%%%%%%%%%%%%%%%%%%%%%%%%%%%%%%%%%%%%%%%%%%%%%%%%%%%%%%% All experiments are designed to be reproducible using publicly available software and documented configurations. The release package will include: \begin{itemize} \item fixed random seeds; \item model hyperparameters; \item training configurations; \item evaluation protocols; \item hardware requirements. \end{itemize} The experimental workflow is summarized as: \begin{equation} Dataset \rightarrow Training \rightarrow Evaluation \rightarrow Analysis. \end{equation} %%%%%%%%%%%%%%%%%%%%%%%%%%%%%%%%%%%%%%%%%%%%%%%%%%%%%%%%%%%%%%%%%%%%%%%%%%%%%%%% \subsection{Software Environment} %%%%%%%%%%%%%%%%%%%%%%%%%%%%%%%%%%%%%%%%%%%%%%%%%%%%%%%%%%%%%%%%%%%%%%%%%%%%%%%% The implementation is based on open-source scientific computing frameworks, including: \begin{itemize} \item Python; \item PyTorch; \item PyTorch Geometric; \item NumPy; \item SciPy; \item Matplotlib. \end{itemize} The complete software environment specification will be provided together with the source code repository. %%%%%%%%%%%%%%%%%%%%%%%%%%%%%%%%%%%%%%%%%%%%%%%%%%%%%%%%%%%%%%%%%%%%%%%%%%%%%%%% \subsection{License} %%%%%%%%%%%%%%%%%%%%%%%%%%%%%%%%%%%%%%%%%%%%%%%%%%%%%%%%%%%%%%%%%%%%%%%%%%%%%%%% The released source code and datasets will be distributed under open-source licenses to encourage scientific collaboration and reuse. The final license information will be specified in the corresponding repository and data archive. %%%%%%%%%%%%%%%%%%%%%%%%%%%%%%%%%%%%%%%%%%%%%%%%%%%%%%%%%%%%%%%%%%%%%%%%%%%%%%%% \subsection{Summary} %%%%%%%%%%%%%%%%%%%%%%%%%%%%%%%%%%%%%%%%%%%%%%%%%%%%%%%%%%%%%%%%%%%%%%%%%%%%%%%% The authors aim to provide a complete and transparent research package including source code, benchmark datasets, and experimental configurations. This release will enable independent verification of PHYSGEN results and support future developments in physics-guided scientific machine learning. \] 
